\section{Introduction}

\subsection{Existing Work}

Sequential commons provide a controlled setting for studying governance when individual incentives can damage a shared resource. The classical common-pool resource literature identifies the same problem in fisheries, irrigation systems, forests, and other renewable resources: open access and weak institutions are associated with overuse even when long-run collective welfare would favour restraint \citep{gordon1954common, scott1955fishery, ostrom1990}. Later work links institutional performance to monitoring, incentives, local participation, and the relation between local and higher-level authority \citep{hilborn2005, gutierrez2011, ostrom1993, ostrom1994, villamayor2020, nagendra2012}. These studies provide the substantive basis for comparing local, central, and mixed governance arrangements in shared-resource settings.

Sequential social-dilemma and multi-agent benchmark work develops this problem in computational environments. These benchmarks make it possible to study repeated interaction, partner composition, learning dynamics, and held-out social settings in a controlled way \citep{leibo2017multiagent, meltingpot2021, socialjax2025}. They are useful for identifying failures that do not appear in single-agent evaluation. A policy can perform well in isolation and still be associated with poor collective outcomes once other agents respond strategically or once the environment rewards short-term extraction.

Scalable oversight studies a related problem from a different direction. Work on scalable oversight asks how a weaker overseer can evaluate, guide, or check a stronger system, and how oversight performance changes as the capability gap widens \citep{bowman2022oversight, engels2025scaling}. The methodological lesson for this study is that oversight should not be treated as perfect by default. Detection, intervention capacity, delay, and cost can be included as evaluation conditions.

Recent work on LLM-generated populations adds a further motivation for studying changing strategic populations. Willis, Du, and Leibo study language-model-generated strategies in a social dilemma, and Willis, Zhao, Du, and Leibo scale this idea to larger populations with cultural evolution \citep{willis2025cooperation, willis2026collective}. Their results report systematic differences in cooperative behaviour across model-generated strategies and poor collective outcomes under some population dynamics. The present paper uses this concern about changing populations to motivate adversarial turnover in a commons benchmark; direct evaluation of LLM strategy populations is left for later work.

\subsection{Research Question}

The research question is: which governance architectures remain robust in sequential commons when stronger adversarial strategies repeatedly enter the population and oversight is imperfect? The study focuses on governance architecture instead of a single learned policy. The comparison covers bottom-up-only, top-down-only, and hybrid governance, with a no-governance condition used as a baseline in the broader evidence chain.

The study evaluates this question in Harvest Commons. Harvest is re-evaluated under three literature-backed scenario archetypes corresponding to regulated fishery, community irrigation, and forest co-management. Oversight is represented through four implementation constraints: imperfect detection, delayed enforcement, limited targeting capacity, and governance budget cost. In operational terms, these constraints act through the government agent. They determine whether an intended cap is detected, when it becomes active, how many agents can be targeted in one step, and what cost is charged when intervention occurs.

The empirical result is a scenario-dependent comparison of governance architecture under explicit oversight frictions. Community-irrigation settings favour hybrid governance under both ideal and constrained oversight. Constrained forest co-management and the higher-pressure regulated-fishery setting favour top-down-only governance. The benchmark reports how architecture rankings vary when adversarial turnover is retained and oversight quality is varied directly.
